\documentclass[12pt,a4paper]{article}
\usepackage[utf8]{inputenc}
\usepackage[spanish,german]{babel}
\usepackage[T1]{fontenc}
\usepackage{geometry}
\geometry{a4paper, margin=2.5cm}
\usepackage{titlesec}
\usepackage{enumitem}
\usepackage{multicol}
\usepackage{csquotes}

\titleformat{\section}
{\Large\bfseries}
{}
{0em}
{}[\titlerule]

\titleformat{\subsection}
{\large\bfseries}
{}
{0em}
{}

\setlength{\parindent}{0pt}
\setlength{\parskip}{0.5em}

\begin{document}

\title{\textbf{DIALOG: En el Transporte Público}}
\author{Spanisch A1/A2}
\date{}
\maketitle

\textbf{Tema:} En el transporte público \\
\textbf{Nivel:} A1-A2 \\
\textbf{Complejidad:} Leicht-Mittel \\
\textbf{Palabras:} 812

\vspace{1em}

\textit{Nota: Se recomienda revisar primero el vocabulario antes de leer el diálogo.}

\vspace{1em}

\section*{Diálogo}

\textbf{Carlos:} Perdón, ¿me puedes ayudar? No sé cómo usar el transporte público aquí.

\textbf{María:} Claro, con mucho gusto. ¿Adónde quieres ir?

\textbf{Carlos:} Quiero ir al centro. ¿Qué bus debo tomar?

\textbf{María:} Toma el bus número 45. Te lleva directamente al centro.

\textbf{Carlos:} ¿Dónde lo compro? ¿Dónde compro el boleto?

\textbf{María:} Puedes comprarlo en la máquina que está ahí. Mira, esa máquina verde.

\textbf{Carlos:} Ah, ya la veo. ¿Cómo la uso?

\textbf{María:} Es fácil. Primero, selecciona tu destino. Luego, inserta el dinero. Finalmente, toma el boleto que te da.

\textbf{Carlos:} Entiendo. ¿Cuánto cuesta?

\textbf{María:} Cuesta dos mil pesos. ¿Tienes monedas?

\textbf{Carlos:} Sí, las tengo. Déjame comprarlo ahora.

\textbf{María:} Perfecto. Mientras lo compras, te explico cómo usarlo.

\textbf{Carlos:} Gracias, eso me ayuda mucho.

\textbf{María:} Cuando subas al bus, muestra el boleto al conductor. Él lo revisará.

\textbf{Carlos:} ¿Y dónde me siento?

\textbf{María:} Puedes sentarte donde quieras, pero primero deja salir a las personas que bajan. No las bloquees.

\textbf{Carlos:} Claro, entiendo. ¿Cómo sé dónde bajar?

\textbf{María:} El bus tiene un sistema de anuncios. Escúchalos. Cuando escuches \enquote{centro}, esa es tu parada.

\textbf{Carlos:} ¿Y si no los escucho bien?

\textbf{María:} Puedes pedirle al conductor que te avise. Dile: \enquote{Por favor, avíseme cuando lleguemos al centro}.

\textbf{Carlos:} Perfecto. ¿Y cómo salgo del bus?

\textbf{María:} Cuando el bus se detenga, presiona el botón rojo que está cerca de la puerta. La puerta se abrirá y puedes salir.

\textbf{Carlos:} ¿Y si hay mucha gente?

\textbf{María:} Di \enquote{con permiso} y las personas te dejarán pasar. No las empujes, solo pide permiso.

\textbf{Carlos:} Entiendo. ¿Y el metro? ¿Cómo funciona?

\textbf{María:} El metro es similar. Primero, compra tu tarjeta en la estación. Luego, cárgala con dinero. Finalmente, úsala para entrar.

\textbf{Carlos:} ¿Dónde la cargo?

\textbf{María:} La cargas en las máquinas que están en cada estación. Son muy fáciles de usar.

\textbf{Carlos:} ¿Y cómo sé qué línea tomar?

\textbf{María:} Mira el mapa que está en la pared. Allí lo verás todo. El mapa te muestra todas las líneas y estaciones.

\textbf{Carlos:} ¿Y si me pierdo?

\textbf{María:} No te preocupes. Pregunta a cualquier persona. La gente es muy amable y te ayudará.

\textbf{Carlos:} ¿Y si no entiendo las indicaciones?

\textbf{María:} Pide ayuda. Di: \enquote{Disculpe, ¿me puede ayudar?} o \enquote{¿Puede explicarme cómo llegar?}

\textbf{Carlos:} Perfecto. Oye, ¿qué pasa si pierdo mi tarjeta?

\textbf{María:} Si la pierdes, debes comprar una nueva. No puedes recuperar el dinero que tenía.

\textbf{Carlos:} Entiendo. ¿Y los horarios? ¿A qué hora pasa el último bus?

\textbf{María:} El último bus pasa a las once de la noche. Después de esa hora, solo hay buses nocturnos, pero son más caros.

\textbf{Carlos:} ¿Y el metro?

\textbf{María:} El metro cierra a medianoche y abre a las cinco de la mañana.

\textbf{Carlos:} Perfecto. Gracias por toda tu ayuda. Ahora entiendo cómo funciona.

\textbf{María:} De nada. Si necesitas más ayuda, no dudes en preguntar. ¡Que tengas un buen viaje!

\textbf{Carlos:} Gracias. ¡Que tengas un buen día!

\newpage

\section*{Vocabulario}

\begin{multicols}{2}
\begin{itemize}[leftmargin=*]
    \item \textbf{transporte público} - öffentlicher Verkehr
    \item \textbf{bus} - Bus
    \item \textbf{metro} - U-Bahn, Metro
    \item \textbf{boleto} - Fahrkarte, Ticket
    \item \textbf{tarjeta} - Karte
    \item \textbf{comprar} - kaufen
    \item \textbf{máquina} - Automat, Maschine
    \item \textbf{insertar} - einwerfen, einstecken
    \item \textbf{dinero} - Geld
    \item \textbf{monedas} - Münzen
    \item \textbf{conductor} - Fahrer, Busfahrer
    \item \textbf{revisar} - überprüfen, kontrollieren
    \item \textbf{subir} - einsteigen, hochsteigen
    \item \textbf{bajar} - aussteigen, heruntersteigen
    \item \textbf{parada} - Haltestelle
    \item \textbf{anuncios} - Durchsagen, Ankündigungen
    \item \textbf{avisar} - benachrichtigen, Bescheid sagen
    \item \textbf{detenerse} - anhalten
    \item \textbf{botón} - Knopf, Taste
    \item \textbf{puerta} - Tür
    \item \textbf{con permiso} - Entschuldigung (um durchzugehen)
    \item \textbf{empujar} - schieben, drängen
    \item \textbf{estación} - Station, Haltestelle
    \item \textbf{cargar} - aufladen (Geld auf Karte)
    \item \textbf{línea} - Linie
    \item \textbf{mapa} - Karte, Plan
    \item \textbf{perderse} - sich verlaufen
    \item \textbf{indicaciones} - Anweisungen, Hinweise
    \item \textbf{horarios} - Fahrpläne, Zeiten
    \item \textbf{buses nocturnos} - Nachtbusse
    \item \textbf{disculpe} - Entschuldigung
    \item \textbf{recuperar} - zurückbekommen
    \item \textbf{medianoche} - Mitternacht
    \item \textbf{bloquear} - blockieren
    \item \textbf{seleccionar} - auswählen
    \item \textbf{destino} - Ziel
\end{itemize}
\end{multicols}

\newpage

\section*{Preguntas de Comprensión}

\begin{enumerate}[leftmargin=*]
    \item ¿Qué quiere hacer Carlos al principio del diálogo?
    \item ¿Qué bus debe tomar Carlos para ir al centro?
    \item ¿Dónde puede comprar Carlos el boleto?
    \item ¿Cuánto cuesta el boleto?
    \item ¿Qué debe hacer Carlos cuando sube al bus?
    \item ¿Qué debe hacer Carlos antes de sentarse?
    \item ¿Cómo sabe Carlos dónde bajar del bus?
    \item ¿Qué puede hacer Carlos si no escucha bien los anuncios?
    \item ¿Cómo sale Carlos del bus?
    \item ¿Qué debe decir Carlos si hay mucha gente al salir?
    \item ¿Cómo funciona el metro según María?
    \item ¿Dónde puede cargar Carlos su tarjeta del metro?
    \item ¿Dónde puede ver Carlos el mapa del metro?
    \item ¿Qué debe hacer Carlos si se pierde?
    \item ¿Qué pasa si Carlos pierde su tarjeta?
    \item ¿A qué hora pasa el último bus?
    \item ¿A qué hora cierra el metro?
    \item ¿Qué son los buses nocturnos?
    \item ¿Qué expresión usa María para decir que Carlos puede preguntar?
    \item ¿Cómo se despide María de Carlos?
\end{enumerate}

\newpage

\section*{Verbo Irregular: SALIR (ausgehen, verlassen)}

\textit{Nota: \enquote{vosotros/as} se usa solo en España, no en español sudamericano. Se incluye aquí para completitud.}

\begin{multicols}{2}

\subsection*{Presente}
\begin{itemize}[leftmargin=*]
    \item yo - salgo
    \item tú - sales
    \item él/ella/usted - sale
    \item nosotros/as - salimos
    \item vosotros/as - salís
    \item ustedes - salen
    \item ellos/ellas - salen
\end{itemize}

\subsection*{Pretérito Indefinido}
\begin{itemize}[leftmargin=*]
    \item yo - salí
    \item tú - saliste
    \item él/ella/usted - salió
    \item nosotros/as - salimos
    \item vosotros/as - salisteis
    \item ustedes - salieron
    \item ellos/ellas - salieron
\end{itemize}

\subsection*{Pretérito Imperfecto}
\begin{itemize}[leftmargin=*]
    \item yo - salía
    \item tú - salías
    \item él/ella/usted - salía
    \item nosotros/as - salíamos
    \item vosotros/as - salíais
    \item ustedes - salían
    \item ellos/ellas - salían
\end{itemize}

\columnbreak

\subsection*{Futuro Simple}
\begin{itemize}[leftmargin=*]
    \item yo - saldré
    \item tú - saldrás
    \item él/ella/usted - saldrá
    \item nosotros/as - saldremos
    \item vosotros/as - saldréis
    \item ustedes - saldrán
    \item ellos/ellas - saldrán
\end{itemize}

\subsection*{Pretérito Perfecto}
\begin{itemize}[leftmargin=*]
    \item yo - he salido
    \item tú - has salido
    \item él/ella/usted - ha salido
    \item nosotros/as - hemos salido
    \item vosotros/as - habéis salido
    \item ustedes - han salido
    \item ellos/ellas - han salido
\end{itemize}

\end{multicols}

\newpage

\section*{Explicación Gramatical}

\textbf{Direkte Objektpronomen und Imperativ}

\textbf{1. Direkte Objektpronomen}

Direkte Objektpronomen ersetzen das direkte Objekt (Akkusativobjekt) im Satz. Sie stehen normalerweise vor dem konjugierten Verb oder werden an den Infinitiv oder das Gerundium angehängt.

Die direkten Objektpronomen sind:
\begin{itemize}[leftmargin=*]
    \item \textbf{me} - mich
    \item \textbf{te} - dich
    \item \textbf{lo} - ihn, es (maskulin)
    \item \textbf{la} - sie, es (feminin)
    \item \textbf{nos} - uns
    \item \textbf{os} - euch (nur in Spanien)
    \item \textbf{los} - sie (maskulin, Plural)
    \item \textbf{las} - sie (feminin, Plural)
\end{itemize}

\textbf{2. Imperativ (Befehlsform)}

Der Imperativ wird verwendet, um Befehle, Anweisungen oder Ratschläge zu geben. Im Spanischen gibt es verschiedene Formen:

\begin{itemize}[leftmargin=*]
    \item \textbf{Tú-Imperativ} (informell, Singular): Wird durch Weglassen des \enquote{-s} von der 3. Person Singular gebildet (bei regelmäßigen Verben) oder hat unregelmäßige Formen.
    \item \textbf{Usted/ustedes-Imperativ} (formell): Wird mit dem Subjuntivo Presente gebildet.
    \item \textbf{Nosotros-Imperativ}: Wird mit dem Subjuntivo Presente gebildet (z.B. \enquote{vamos}).
\end{itemize}

\textbf{3. Direkte Objektpronomen mit Imperativ}

Bei positiven Imperativformen werden die Objektpronomen an das Verb angehängt. Bei negativen Imperativformen stehen sie davor.

\textbf{Ejemplos del texto:}

\textbf{Direkte Objektpronomen:}
\begin{itemize}[leftmargin=*]
    \item \enquote{¿Dónde \textbf{lo} compro?} -- Wo kaufe ich \textbf{es}? (\textbf{lo} = el boleto)
    \item \enquote{¿Dónde \textbf{la} cargo?} -- Wo lade ich \textbf{sie} auf? (\textbf{la} = la tarjeta)
    \item \enquote{Él \textbf{lo} revisará} -- Er wird \textbf{es} überprüfen (\textbf{lo} = el boleto)
    \item \enquote{¿Tienes monedas? Sí, \textbf{las} tengo} -- Hast du Münzen? Ja, ich habe \textbf{sie} (\textbf{las} = las monedas)
    \item \enquote{¿Me puedes ayudar?} -- Kannst du \textbf{mir} helfen? (\textbf{me} = mir, indirektes Objektpronomen, aber ähnliche Struktur)
    \item \enquote{No \textbf{las} bloquees} -- Blockiere \textbf{sie} nicht (\textbf{las} = las personas)
    \item \enquote{Si \textbf{la} pierdes} -- Wenn du \textbf{sie} verlierst (\textbf{la} = la tarjeta)
\end{itemize}

\textbf{Imperativ:}
\begin{itemize}[leftmargin=*]
    \item \enquote{\textbf{Toma} el bus número 45} -- Nimm den Bus Nummer 45 (tú-Imperativ von \enquote{tomar})
    \item \enquote{\textbf{Mira}, esa máquina verde} -- Schau, diese grüne Maschine (tú-Imperativ von \enquote{mirar})
    \item \enquote{\textbf{Di} \enquote{con permiso}} -- Sag \enquote{Entschuldigung} (tú-Imperativ von \enquote{decir})
    \item \enquote{No \textbf{las} empujes} -- Dränge \textbf{sie} nicht (negativer tú-Imperativ von \enquote{empujar})
    \item \enquote{\textbf{Pide} ayuda} -- Bitte um Hilfe (tú-Imperativ von \enquote{pedir})
    \item \enquote{\textbf{Escúchalos}} -- Hör \textbf{sie} an (tú-Imperativ + angehängtes Objektpronomen \enquote{los})
    \item \enquote{\textbf{Avíseme} cuando lleguemos} -- Benachrichtigen Sie \textbf{mich}, wenn wir ankommen (usted-Imperativ von \enquote{avisar} + \enquote{me})
\end{itemize}

\textbf{Wichtige Hinweise:}
\begin{itemize}[leftmargin=*]
    \item Bei positiven Imperativformen werden die Objektpronomen an das Verb angehängt: \enquote{¡Cómpralo!} (Kauf es!)
    \item Bei negativen Imperativformen stehen die Objektpronomen vor dem Verb: \enquote{¡No lo compres!} (Kauf es nicht!)
    \item Wenn ein Objektpronomen angehängt wird, kann ein Akzent hinzugefügt werden müssen, um die Betonung beizubehalten: \enquote{compra} → \enquote{cómpralo}
    \item In südamerikanischem Spanisch wird \enquote{ustedes} statt \enquote{vosotros} verwendet, daher ist der ustedes-Imperativ häufiger.
\end{itemize}

\end{document}
